% pdflatex-tate-math-movarg.tex -- P6: moving argument 内の数式の保護
% $ は cat3 のままなので \protected@edef/\write/\mark で捕捉機構は展開
% されない。判定（runner: 2パス + ログ検査）:
%   SECTION : 見出し移動引数内の $..$ が aux を破壊しない（2パス成立）
%   MARK    : \mark 中の数式でエラーにならない
%   BOOKMARK: hyperref しおり文字列の数式フォールバックが機能（.out に
%             数式由来のASCII列が残る; hyperref 既存フォールバック）
%   REFS    : 式番号 \label/\eqref が解決（2パス目で未解決警告なし）
%   BODY    : 本文の行内数式も捕捉経路で成立
% NOTE: tate + \tableofcontents のドットリーダーは本 P6 とは無関係の
% 既存不良（HEAD でも発生・数式なしでも2パス目で "Leaders not
% followed by proper glue"。TOC 実測は [tate]{jlreq} の出力ルーチン
% スライス (P7) で回収予定のため、ここでは目次の組版は対象外とし
% 移動引数の保護のみを検証する）。
\documentclass[pdflatex,tate,paper=a5]{jlreq}
\usepackage{amsmath}
\usepackage[unicode]{hyperref}
\pagestyle{empty}
\makeatletter
\begin{document}
\section{移動引数の試験 $a+b$}\label{sec:mov}
本文: 式~\eqref{eq:mv} を参照。節~\ref{sec:mov}。
\begin{equation}
  x + y = z \label{eq:mv}
\end{equation}
\markboth{柱の試験 $m$}{柱の試験 $m$}
漢字$a+b$漢字
\typeout{TATE-PASS: SECTION (heading with math in moving arg)}%
\typeout{TATE-PASS: MARK (mark with math ok)}%
\typeout{TATE-PASS: BOOKMARK (hyperref math fallback; .out checked by runner)}%
\typeout{TATE-PASS: BODY (inline math in body text)}%
\end{document}
